% =============================================================================
% Notation and Conventions -- frontmatter, ungrouped by section, referenced
% throughout the body and appendices. A symbol used with a different LOCAL
% meaning inside one short, self-contained derivation (e.g. Appendix B's own
% $a_\ell, p_\ell$ for a function and its derivative, before the construction
% is restated in the $c^{(j)}$ notation below) is resolved back to this table
% before it recurs elsewhere, and is not itself listed here.
% =============================================================================

A quick reference for symbols used across the thesis; each is properly
introduced at its first use in the text below.

\bigskip
\noindent\textbf{Discretisation}
\begin{center}
\begin{tabular}{@{}l p{0.82\textwidth}@{}}
\toprule
Symbol & Meaning \\
\midrule
$n$ & Number of qubits in the Chebyshev coefficient register (the resolution). \\
$N = 2^n$ & Chebyshev truncation degree; dimension of the single-field Hilbert space. \\
\bottomrule
\end{tabular}
\end{center}

\bigskip
\noindent\textbf{Differential-equation structure}
\begin{center}
\begin{tabular}{@{}l p{0.82\textwidth}@{}}
\toprule
Symbol & Meaning \\
\midrule
$k$ & Order of an \gls{ode}. \\
$p$ & Polynomial (nonlinear) degree of the \gls{ode} (\S4). \\
\bottomrule
\end{tabular}
\end{center}

\bigskip
\noindent\textbf{Chebyshev and descriptor operators}
\begin{center}
\begin{tabular}{@{}l p{0.82\textwidth}@{}}
\toprule
Symbol & Meaning \\
\midrule
$T_\ell$ & Chebyshev polynomial of the first kind, degree $\ell$. \\
$G$ & Dense $N \times N$ Chebyshev differentiation matrix (the collocation route's operator; Wu et al.\ \cite{wu2025pihm}). \\
$\tilde{B}, \tilde{D}$ & Banded factors of $G$ (\Cref{app:cheb}). \\
$M$ & Coefficient-multiplication operator in Chebyshev coefficient space (\Cref{app:cheb}). \\
$w$ & The stacked descriptor state, $w = (c^{(0)}; \dots; c^{(k)})$. \\
$A$ & The residual operator of a linear-system reformulation of the \gls{ode} (dense in the collocation route, banded in the descriptor route). \\
\bottomrule
\end{tabular}
\end{center}

\bigskip
\noindent\textbf{Physical and complexity quantities}
\begin{center}
\begin{tabular}{@{}l p{0.82\textwidth}@{}}
\toprule
Symbol & Meaning \\
\midrule
$H$ & The physics-informed effective Hamiltonian (Wu et al.\ \cite{wu2025pihm}). \\
$\alpha$ & Block-encoding subnormalisation factor. \\
$\Delta$ & Spectral gap of $H$ between its ground state and first excited state. \\
$\epsilon_G$ & Block-encoding / truncation error associated with representing $G$. \\
\bottomrule
\end{tabular}
\end{center}
